% !TEX program = pdflatex
%=====================================================================
%  Criptografía y Seguridad - ITBA
%  Clase 10b — Probando la seguridad: Penetration testing
%  Apunte integral (teoría; esta clase no tiene práctica asociada)
%=====================================================================
\documentclass[11pt,a4paper]{article}

%---------------------------------------------------------------------
%  Paquetes
%---------------------------------------------------------------------
\usepackage[utf8]{inputenc}
\usepackage[T1]{fontenc}
\usepackage[spanish,es-tabla,es-noshorthands]{babel}
\usepackage{lmodern}
\usepackage{microtype}

\usepackage{amsmath,amssymb,amsthm}
\usepackage{mathtools}
\usepackage{geometry}
\geometry{a4paper,margin=2.4cm,headheight=22pt}

\usepackage{xcolor}
\usepackage{graphicx}
\usepackage{enumitem}
\usepackage{booktabs}
\usepackage{array}
\usepackage{tcolorbox}
\tcbuselibrary{skins,breakable}
\usepackage{fancyhdr}
\usepackage{titlesec}
\usepackage{hyperref}
\usepackage{listings}
\usepackage{caption}

\usepackage{tikz}
\usetikzlibrary{shapes.geometric,shapes.misc,shapes.symbols,arrows.meta,positioning,calc,
                fit,backgrounds,decorations.pathmorphing,shadows.blur,chains}

%---------------------------------------------------------------------
%  Paleta de color (inspirada en las diapositivas de la cátedra)
%---------------------------------------------------------------------
\definecolor{itbaCyan}{HTML}{6EC6D9}
\definecolor{itbaCyanDark}{HTML}{2E8B9E}
\definecolor{itbaBlue}{HTML}{AEDCEA}
\definecolor{boxBlue}{HTML}{BBD4F0}
\definecolor{slideGray}{HTML}{ECECEC}
\definecolor{practGreen}{HTML}{4F8A3D}
\definecolor{practGreenL}{HTML}{E2EFD9}
\definecolor{attackRed}{HTML}{C0392B}
\definecolor{codeBg}{HTML}{F5F5F5}

%---------------------------------------------------------------------
%  Hipervínculos
%---------------------------------------------------------------------
\hypersetup{
  colorlinks=true,
  linkcolor=itbaCyanDark,
  urlcolor=itbaCyanDark,
  citecolor=itbaCyanDark,
  pdftitle={Clase 10b - Penetration Testing},
  pdfauthor={Criptografia y Seguridad - ITBA}
}

%---------------------------------------------------------------------
%  Encabezados y pies
%---------------------------------------------------------------------
\pagestyle{fancy}
\fancyhf{}
\lhead{\small\textsc{Criptografía y Seguridad — ITBA}}
\rhead{\small Clase 10b · Penetration testing}
\cfoot{\small\thepage}
\renewcommand{\headrulewidth}{0.4pt}
\renewcommand{\headrule}{\hbox to\headwidth{\color{itbaCyanDark}\leaders\hrule height \headrulewidth\hfill}}

%---------------------------------------------------------------------
%  Títulos de sección con barra de color (estilo diapositiva)
%---------------------------------------------------------------------
\titleformat{\section}
  {\normalfont\Large\bfseries\color{itbaCyanDark}}
  {\thesection}{0.6em}{}
  [\vspace{-0.4em}{\color{itbaCyan}\titlerule[2pt]}]
\titleformat{\subsection}
  {\normalfont\large\bfseries\color{itbaCyanDark}}
  {\thesubsection}{0.6em}{}
\titleformat{\subsubsection}
  {\normalfont\normalsize\bfseries\color{black}}
  {\thesubsubsection}{0.6em}{}

%---------------------------------------------------------------------
%  Cajas reutilizables
%---------------------------------------------------------------------
% Caja "diapositiva" (recuadro gris de las slides)
\newtcolorbox{slidebox}[1][]{
  enhanced, colback=slideGray, colframe=black!55, boxrule=0.5pt,
  arc=1pt, left=6pt,right=6pt,top=4pt,bottom=4pt, #1}

% Caja de definición
\newtcolorbox{defbox}[1]{
  enhanced, breakable, colback=itbaBlue!25, colframe=itbaCyanDark,
  boxrule=1pt, arc=2pt, left=8pt,right=8pt,top=6pt,bottom=6pt,
  fonttitle=\bfseries, coltitle=white, title={#1}}

% Caja de nota / aclaración del profe
\newtcolorbox{notebox}[1]{
  enhanced, breakable, colback=practGreenL, colframe=practGreen,
  boxrule=1pt, arc=2pt, left=8pt,right=8pt,top=6pt,bottom=6pt,
  fonttitle=\bfseries, coltitle=white, title={#1}}

% Caja de atención / ataque
\newtcolorbox{warnbox}[1]{
  enhanced, breakable, colback=attackRed!8, colframe=attackRed,
  boxrule=1pt, arc=2pt, left=8pt,right=8pt,top=6pt,bottom=6pt,
  fonttitle=\bfseries, coltitle=white, title={#1}}

%---------------------------------------------------------------------
%  Macros
%---------------------------------------------------------------------
\newcommand{\factor}[1]{\textcolor{itbaCyanDark}{\textbf{#1}}}
\newcommand{\prof}[1]{\medskip\noindent{\small\itshape``#1''}\\[-1pt]
  \hspace*{1em}{\footnotesize\textcolor{black!55}{--- Prof.\ Leandro Suárez, en clase}}\medskip}

%=====================================================================
\begin{document}
%=====================================================================

%---------------------------------------------------------------------
%  Portada
%---------------------------------------------------------------------
\begin{titlepage}
\centering
\vspace*{1.2cm}

\begin{tikzpicture}
  \node[rounded corners=2pt, fill=itbaCyan, text=black, inner sep=14pt,
        minimum width=0.8\textwidth, align=center] (t)
        {\Huge\bfseries Criptografía y Seguridad};
\end{tikzpicture}

\vspace{0.6cm}
{\LARGE Probando la seguridad}\\[0.25cm]
{\Huge\bfseries\color{itbaCyanDark} Penetration testing}

\vspace{0.9cm}

% Ilustración: bóveda abierta (recreación del icono de portada de la slide)
\begin{tikzpicture}[scale=1.0]
  % marco / pared
  \fill[black!10] (-3.6,-2.6) rectangle (3.6,3.0);
  \draw[line width=2pt, black!55] (-3.6,-2.6) rectangle (3.6,3.0);
  % interior iluminado de la bóveda (puerta abierta)
  \fill[itbaCyan!25] (-0.9,-2.2) rectangle (0.9,2.6);
  \foreach \y in {-1.6,-0.8,0.0,0.8,1.6}{
    \draw[black!35,line width=0.6pt] (-0.8,\y) rectangle (0.8,\y+0.5);
  }
  % hoja izquierda de la puerta (abierta hacia afuera)
  \begin{scope}
  \fill[itbaCyanDark!70] (-3.0,-2.2) -- (-1.0,-2.0) -- (-1.0,2.4) -- (-3.0,2.6) -- cycle;
  \draw[line width=2pt, black!70] (-3.0,-2.2) -- (-1.0,-2.0) -- (-1.0,2.4) -- (-3.0,2.6) -- cycle;
  \draw[line width=1pt, itbaCyan] (-2.7,-1.8) -- (-1.25,-1.65) -- (-1.25,2.05) -- (-2.7,2.25) -- cycle;
  \draw[line width=2.5pt, black!70] (-2.0,0.2) circle (0.34);
  \foreach \a in {0,60,...,300}{ \draw[line width=1.5pt,black!70] (-2.0,0.2)--++(\a:0.5);}
  \fill[itbaCyan] (-2.0,0.2) circle (0.12);
  \end{scope}
  % hoja derecha de la puerta (abierta hacia afuera)
  \begin{scope}
  \fill[itbaCyanDark!70] (3.0,-2.2) -- (1.0,-2.0) -- (1.0,2.4) -- (3.0,2.6) -- cycle;
  \draw[line width=2pt, black!70] (3.0,-2.2) -- (1.0,-2.0) -- (1.0,2.4) -- (3.0,2.6) -- cycle;
  \draw[line width=1pt, itbaCyan] (2.7,-1.8) -- (1.25,-1.65) -- (1.25,2.05) -- (2.7,2.25) -- cycle;
  \draw[line width=2.5pt, black!70] (2.0,0.2) circle (0.34);
  \foreach \a in {0,60,...,300}{ \draw[line width=1.5pt,black!70] (2.0,0.2)--++(\a:0.5);}
  \fill[itbaCyan] (2.0,0.2) circle (0.12);
  \end{scope}
  % piso
  \draw[line width=2pt,black!55] (-3.6,-2.6) -- (3.6,-2.6);
\end{tikzpicture}

\vfill
{\large Apunte integral — Teoría}\\[0.2cm]
{\normalsize Clase 10b \quad·\quad Capítulo 23 (secciones 1--2), \emph{Computer Security: Art and Science} (M.~Bishop)}\\[0.4cm]
{\small Instituto Tecnológico de Buenos Aires (ITBA)}

\vspace{0.6cm}
{\footnotesize\itshape Documento elaborado a partir de la transcripción de la clase
teórica (Prof.\ Leandro Suárez) y de las diapositivas oficiales de la cátedra.
Esta clase no tiene clase práctica asociada.}
\end{titlepage}

\tableofcontents
\newpage

%=====================================================================
\section{Introducción y contexto}
%=====================================================================

Esta clase responde a una pregunta de fondo: \textbf{¿cómo verificamos que un sistema es
seguro?} Hay varias formas de corroborar el funcionamiento correcto de un sistema; la clase
contrapone dos de ellas —la \textbf{verificación formal} y la \textbf{prueba de penetración}
(\emph{penetration testing} o \emph{pentesting})— y luego profundiza en esta última: su
metodología, sus etapas y dos ejemplos concretos.

\begin{notebox}{Encuadre del profesor}
El \emph{pentesting} había quedado pendiente de la clase anterior, donde sólo se habían
``vislumbrado'' los conceptos de \emph{red team}, \emph{breach \& attack simulation} y las
etapas de un pentest. El profesor aclara que \textbf{todo esto entra en el parcial}:
\prof{Igual, a ver, esto de la segunda parte entra todo en el parcial; así que es lo mismo
que lo veamos hoy o la semana pasada.}
\end{notebox}

\begin{warnbox}{La gran conclusión que atraviesa toda la clase}
El \emph{testing} puede mostrar la \textbf{presencia} de bugs/vulnerabilidades, pero
\textbf{nunca prueba su ausencia} (idea atribuida en clase a Dijkstra). El pentesting
\emph{parte de la premisa de que ya existen vulnerabilidades} e intenta demostrarlas; que
no encuentre nada \textbf{no} significa que el sistema sea seguro, sólo que no se comprometió
\emph{bajo las condiciones de esa prueba}.
\end{warnbox}

%=====================================================================
\section{Verificación formal}
%=====================================================================

\begin{defbox}{Verificación formal}
Verificación \textbf{matemática} de que un sistema cumple con ciertas restricciones
(las que nosotros le imponemos). Se modela el sistema mediante:
\begin{itemize}[leftmargin=1.4em,itemsep=2pt]
  \item \textbf{Precondiciones}: una hipótesis sobre el \emph{estado} del sistema
        (p.\ ej.\ ``el sistema está en un estado seguro'').
  \item \textbf{Poscondiciones}: el resultado de aplicar las operaciones del sistema a un
        \emph{input dado}.
\end{itemize}
\textbf{Requerimiento:} las poscondiciones deben cumplir las restricciones impuestas.
\end{defbox}

Es el clásico esquema de prueba matemática: se parte de una hipótesis (precondición), se
ejecuta la verificación y se corrobora que las poscondiciones cumplen lo que se pedía. Está
orientada a \emph{probar matemáticamente que un sistema hace lo que dice hacer}.

\subsection*{Diagrama del modelo}

La siguiente figura recrea el esquema de la diapositiva: las \emph{precondiciones} y las
\emph{entradas} alimentan una secuencia de operaciones $\mathrm{Op}_1,\dots,\mathrm{Op}_n$
del sistema, que produce las \emph{poscondiciones}.

\begin{center}
\begin{tikzpicture}[
   node distance=6mm,
   box/.style={draw=black!70, fill=itbaBlue!45, minimum width=2.6cm, minimum height=0.85cm, align=center, font=\small},
   op/.style={draw=black!70, fill=itbaCyan!50, minimum width=0.9cm, minimum height=0.9cm, align=center, font=\footnotesize},
   >=Stealth]
  \node[box] (pre) {Precondiciones};
  \node[box, below=of pre] (in) {Entradas};
  % bloque de operaciones
  \node[draw=black!70, fill=itbaCyan!25, minimum width=3.2cm, minimum height=2.1cm] (sys) at (5.0,-0.65) {};
  \node[op] (op1) at (4.1,-0.65) {$\mathrm{Op}_1$};
  \node at (5.0,-0.65) {$\cdots$};
  \node[op] (opn) at (5.9,-0.65) {$\mathrm{Op}_n$};
  \node[box] (post) at (10.0,-0.65) {Poscondiciones};
  % flechas
  \draw[->] (pre.east) -- ($(sys.west)+(0,0.35)$);
  \draw[->] (in.east)  -- ($(sys.west)+(0,-0.35)$);
  \draw[<-] (post.west) -- (sys.east);
\end{tikzpicture}
\end{center}

\begin{notebox}{¿Por qué casi no se usa en la práctica?}
La verificación formal es \textbf{muy potente pero impráctica} para sistemas reales:
\begin{itemize}[leftmargin=1.4em,itemsep=1pt]
  \item Es \textbf{extremadamente compleja}: con sólo un par de variables y operaciones ya
        cuesta seguirla; hay que mapear todas las posibilidades, precondiciones e
        interacciones entre componentes.
  \item Es \textbf{muy costosa} en tiempo: a veces lleva más tiempo verificar el sistema que
        construirlo de nuevo.
  \item Los sistemas \textbf{evolucionan muy rápido}: antes de terminar la verificación, el
        sistema ya cambió y hay que rehacerla.
\end{itemize}
El profesor la compara con las viejas ``pruebas de escritorio'': probar una función con
distintos inputs (incluido el caso \texttt{null}) es realizable; probar \emph{un sistema
completo} de esta forma se vuelve impracticable.
\end{notebox}

%=====================================================================
\section{Prueba de penetración (\textit{Penetration testing})}
%=====================================================================

\begin{defbox}{Prueba de penetración}
\textbf{Prueba} (no demostración matemática) para verificar que un sistema cumple con ciertas
restricciones. Se modela con el mismo lenguaje de pre/poscondiciones que la verificación
formal, pero con un giro clave:
\begin{itemize}[leftmargin=1.4em,itemsep=2pt]
  \item \textbf{Precondiciones}: hipótesis sobre el estado del sistema \textbf{y la
        existencia de una vulnerabilidad}. \emph{Se asume de entrada que la vulnerabilidad
        existe.}
  \item \textbf{Poscondiciones}: \textbf{sistema comprometido}.
  \item \textbf{Ejecución}: aplicar pruebas para intentar \emph{mover al sistema del estado
        inicial al estado comprometido}.
\end{itemize}
\end{defbox}

Si la ejecución logra ese movimiento, se concluye que el sistema \emph{está comprometido}
bajo esas condiciones. Es una forma \textbf{mucho más económica, simple y eficiente} de
probar un sistema que la verificación formal, y es la que usan las empresas (grandes o
PyMEs). Permite, además, enfocar primero los componentes más críticos de un sistema que
evoluciona rápido.

\begin{warnbox}{No encontrar nada $\neq$ ser seguro}
Que la poscondición ``no se cumpla'' (no encontré la vulnerabilidad) \textbf{no significa}
que no haya fallas: puede haber otro vector de ataque o vectores que esa prueba no exploró.
La prueba de penetración \emph{no prueba la ausencia} de vulnerabilidades.
\end{warnbox}

%=====================================================================
\section{Similitudes y diferencias}
%=====================================================================

\begin{center}
\renewcommand{\arraystretch}{1.35}
\begin{tabular}{>{\raggedright\arraybackslash}p{4.0cm} >{\raggedright\arraybackslash}p{5.0cm} >{\raggedright\arraybackslash}p{5.0cm}}
\toprule
 & \textbf{Prueba de penetración} & \textbf{Verificación formal} \\
\midrule
\textbf{Prueba la existencia de vulnerabilidades} & Sí & Sí \\
\textbf{Prueba la ausencia de vulnerabilidades} & \textcolor{attackRed}{\textbf{NO}} & Sí, pero debe incluir \textbf{TODOS} los factores externos \\
\textbf{Costo / practicidad} & Económica, eficiente, práctica & Muy costosa, a menudo impráctica \\
\bottomrule
\end{tabular}
\end{center}

\medskip
\noindent Ambas técnicas \textbf{prueban la existencia} de vulnerabilidades. La diferencia
está en la \emph{ausencia}: la verificación formal sí puede probarla, pero a condición de
incluir \emph{todos} los factores posibles, incluidos los externos. La prueba de penetración
no puede.

\begin{notebox}{Lo que se ignora en la práctica}
En la práctica, incluso una verificación formal prueba la ausencia de vulnerabilidades en
\textbf{un algoritmo, programa o ambiente determinado}, y \textbf{no en un sistema completo}.
Se ignoran cuestiones como la \textbf{instalación} y el \textbf{uso}, el sistema operativo
sobre el que corre, el estado y el lugar donde está instalado, etc. Esto conecta con la frase
de Dijkstra repetida en clase: el testing muestra la presencia de bugs, no su ausencia.
\end{notebox}

%=====================================================================
\section{Objetivos del pentesting}
%=====================================================================

El objetivo central es \textbf{probar la eficacia de los controles de seguridad de un
sistema}. Esto se logra:

\begin{itemize}[leftmargin=1.6em,itemsep=3pt]
  \item \textbf{Intentando violar la política de seguridad} ya existente. Por eso suele haber
        un \emph{equipo designado} que conoce de seguridad y entiende las implicancias.
  \item \textbf{Ejecutando técnicas similares a las de un atacante}: simular ser un atacante.
        Esto se conoce como \factor{ethical hacking}.
  \item De forma \textbf{análoga a las pruebas manuales} de un sistema: alguien que no
        participó de la construcción intenta romperlo (pasar inputs inválidos, ver cómo
        reacciona). \emph{No reemplaza un buen diseño e implementación}.
  \item Probando \textbf{al sistema como un todo}, no sólo sus aspectos técnicos.
\end{itemize}

\begin{warnbox}{Ética del hacking: comportamiento del pentester}
Al simular ser un atacante, el pentester debe adoptar un \textbf{comportamiento ético}:
\textbf{no} robar datos de clientes que pudiera quedar expuestos, \textbf{no} hacer compras
en nombre de otra persona, etc. Es clave cuidar la \emph{credibilidad} de estas pruebas. Es
también recomendable que las haga \emph{alguien que no intervino en la construcción} del
sistema (análogo al rol de QA: el desarrollador prueba, hace regresiones y pruebas unitarias;
luego otro intenta romperlo).
\end{warnbox}

\begin{notebox}{¿Quién hace los pentests y por qué importan?}
El profesor remarca que esta metodología es \textbf{base de lo que hoy se usa para testear},
fundamental si se aspira a trabajar en equipos de \emph{Red Team}, \emph{AppSec} o
\emph{InfraSec}. Para cualquier sistema mediano o grande hoy se \textbf{exige} algún tipo de
pentesting, sobre todo en los componentes \emph{core}, y que no tengan vulnerabilidades
críticas o altas (criticidad definida por los \textbf{CVE}). Por eso es importante la
organización \textbf{MITRE}. A veces se contrata un proveedor; si la empresa es
suficientemente madura y grande, se hace internamente con gente experta.
\end{notebox}

%=====================================================================
\section{Metodología (informal)}
%=====================================================================

Es la metodología vista ``por encima'' la clase anterior. A grandes rasgos, sus pasos son:

\begin{enumerate}[leftmargin=1.8em,itemsep=3pt]
  \item \textbf{Determinar y cuantificar objetivos.} Preparación: qué vamos a hacer, qué
        alcance tiene, con qué tiempo contamos, cuál es el propósito.
        \emph{Ejemplo: obtener información de clientes.}
  \item \textbf{Encontrar cierta cantidad de vulnerabilidades o buscarlas durante un período
        de tiempo.} El estudio se conduce \textbf{desde el punto de vista de un atacante}.
        Aquí se arma el \emph{plan de ataque}: qué herramientas se usan, dónde se ubican, qué
        equipo está a cargo, quiénes participan y quiénes están al tanto.
  \item \textbf{Estudiar y categorizar los hallazgos.} \emph{El éxito de una prueba de
        penetración proviene del análisis de los hallazgos}: permite detectar problemas
        recurrentes y corregir sistemáticamente \textbf{familias} de problemas.
\end{enumerate}

\begin{notebox}{Aclaraciones de la clase}
\begin{itemize}[leftmargin=1.4em,itemsep=2pt]
  \item \textbf{Lower line / aviso:} muchas veces \emph{no} se avisa a los equipos de
        desarrollo u operativos que va a ejecutarse un pentest (para que la prueba sea
        realista).
  \item \textbf{Condición de éxito:} el objetivo tiene una condición que determina si la
        prueba tuvo éxito o no. Se ataca hasta lograr el objetivo o hasta que \textbf{se vence
        el tiempo} definido; ahí termina la prueba.
  \item \textbf{Es cíclico:} como casi todos los procesos de seguridad, los hallazgos se
        reincorporan al ciclo y se encadenan con otras pruebas. Una vulnerabilidad suele tener
        un \emph{patrón de comportamiento}: aparece en distintos componentes (porque se
        desarrollaron igual o usan la misma dependencia vulnerable).
\end{itemize}
\end{notebox}

%=====================================================================
\section{Metodología de Hipótesis de Falla}
%=====================================================================

\begin{warnbox}{¡Importante para el parcial!}
\prof{Esta metodología es importante, es súper importante que la entiendan y la sepan: es la
base de lo que hoy se usa para testear.}
Es la metodología en la que se basa el pentest. Si bien en la práctica no se respeta al
detalle, es la \textbf{estructura} de cómo debería funcionar un pentest a grandes rasgos.
\end{warnbox}

Consta de \textbf{cinco etapas}, la última opcional:

\begin{center}
\begin{tikzpicture}[
  >=Stealth, node distance=4mm,
  stage/.style={draw=itbaCyanDark, fill=itbaBlue!30, rounded corners=3pt,
                minimum width=3.0cm, minimum height=1.0cm, align=center, font=\small},
  opt/.style={draw=itbaCyanDark, fill=slideGray, rounded corners=3pt,
              minimum width=3.0cm, minimum height=1.0cm, align=center, font=\small},
  num/.style={circle, draw=itbaCyanDark, fill=white, font=\bfseries\footnotesize, inner sep=1.5pt}]
  \node[stage] (s1) {Recolección\\de información};
  \node[stage, right=of s1] (s2) {Hipótesis};
  \node[stage, right=of s2] (s3) {Prueba};
  \node[stage, below=8mm of s3] (s4) {Generalización};
  \node[opt,   left=of s4] (s5) {Eliminación\\\footnotesize(opcional)};
  \node[num] at (s1.north west) {1};
  \node[num] at (s2.north west) {2};
  \node[num] at (s3.north west) {3};
  \node[num] at (s4.north west) {4};
  \node[num] at (s5.north west) {5};
  \draw[->] (s1) -- (s2);
  \draw[->] (s2) -- (s3);
  \draw[->] (s3) -- (s4);
  \draw[->] (s4) -- (s5);
  % realimentación (pivoting / ciclo)
  \draw[->, dashed, attackRed] (s4.south) .. controls +(0,-0.8) and +(0,-1.2) .. (s1.south)
        node[midway, below, font=\scriptsize\itshape, text=attackRed] {realimentación (pivoting / ciclo)};
\end{tikzpicture}
\end{center}

\begin{enumerate}[leftmargin=1.8em,itemsep=3pt]
  \item \textbf{Recolección de información} --- para conocer el sistema y su funcionamiento.
  \item \textbf{Hipótesis} --- asumir la existencia de vulnerabilidades concretas.
  \item \textbf{Prueba} --- probar las vulnerabilidades.
  \item \textbf{Generalización} --- generalizar patrones y hallar otras vulnerabilidades del
        mismo tipo.
  \item \textbf{Eliminación (opcional)} --- determinar los pasos necesarios para eliminarlas.
\end{enumerate}

\subsection{Etapa 1 — Recolección de información}

El objetivo es \textbf{componer un modelo del sistema y sus partes}, interiorizarse con él
para saber ``dónde poner el ojo''.

\begin{itemize}[leftmargin=1.6em,itemsep=3pt]
  \item \textbf{Buscar discrepancias.} Cuando el sistema evoluciona y la documentación no
        refleja esa evolución, aparece una posible \emph{puerta de entrada}: la documentación
        dice que se comporta de una forma pero ya no lo hace.
  \item \textbf{Revisar interfaces}: conexiones con sistemas externos, proveedores, otros
        sistemas, bases de datos.
  \item \textbf{Conocer bien el sistema} usando documentación de diseño y manuales (si están
        disponibles). A mayor documentación, mejor \emph{insight}.
        \emph{Especialmente, buscar secciones poco especificadas o ambiguas}: ahí pudo
        tomarse alguna decisión cuestionable $\Rightarrow$ posible puerta de entrada.
  \item \textbf{Ver manejo de privilegios y tipos de cuentas}: enumerar usuarios; buscar
        cuentas con \emph{superadmin}/supervisor (las de mayor poder, generalmente sin
        límites), y \textbf{cuentas huérfanas o con exceso de permisos}.
  \item \textbf{El ambiente}: conseguir nombres de usuarios y servidores, en qué servidores
        corre, la \textbf{estructura de red} del sistema y su configuración.
  \item \textbf{Vulnerabilidades conocidas genéricas} que puedan ser compatibles con la
        tecnología que usa el sistema.
\end{itemize}

\subsection{Punto de partida}

Identifica la \textbf{información inicial con la que cuenta el equipo} (el supuesto atacante).
Se distinguen \textbf{niveles incrementales}:

\begin{center}
\begin{tikzpicture}[
  >=Stealth,
  lvl/.style={draw=black!60, rounded corners=2pt, minimum width=8.2cm, minimum height=0.9cm,
              align=center, font=\small}]
  \node[lvl, fill=itbaBlue!20] (l1) {\textbf{Nivel 1:} Atacante externo \emph{sin} conocimiento del sistema};
  \node[lvl, fill=itbaBlue!35, below=2mm of l1] (l2) {\textbf{Nivel 2:} Atacante externo \emph{con} conocimiento del sistema};
  \node[lvl, fill=itbaBlue!50, below=2mm of l2] (l3) {\textbf{Nivel 3:} Atacante \emph{con acceso} al sistema};
  \draw[->, line width=1pt, itbaCyanDark] (8.6,-0.0) ++(0,0.0) -- ++(0,-2.0)
        node[midway, right, font=\scriptsize, text=itbaCyanDark, align=left] {más\\información};
\end{tikzpicture}
\end{center}

Dependiendo del tipo de prueba, algunos niveles son \textbf{irrelevantes}:
\begin{itemize}[leftmargin=1.6em,itemsep=2pt]
  \item Durante el \textbf{diseño}, se ignora el \emph{primer} nivel.
  \item En sistemas con \textbf{registro abierto}, se ignora el \emph{segundo} nivel.
\end{itemize}

\subsection{Etapa 2 — Hipótesis}

Aquí se usa toda la información recabada para \textbf{buscar posibles vulnerabilidades},
tratando de encontrar combinaciones que expongan una vulnerabilidad.

\subsubsection*{Examinando políticas y procedimientos}
\begin{itemize}[leftmargin=1.6em,itemsep=2pt]
  \item Buscar \textbf{inconsistencias} (incluidas las inconsistencias entre \emph{políticas} y
        \emph{mecanismos} que las implementan).
  \item Los \textbf{procedimientos pueden no cumplirse} en la práctica.
\end{itemize}

\subsubsection*{Examinando implementaciones}
\begin{itemize}[leftmargin=1.6em,itemsep=2pt]
  \item Utilizar \textbf{modelos de vulnerabilidades}.
  \item Revisar vulnerabilidades comunes/conocidas (\factor{vulnerability scanning},
        p.\ ej.\ \texttt{portmapper}).
  \item Usar manuales para \textbf{exceder límites}, \textbf{omitir pasos} de las secuencias,
        etc.
\end{itemize}

\subsubsection*{Vectores de hipótesis específicos (slide de Hipótesis)}
\begin{itemize}[leftmargin=1.6em,itemsep=3pt]
  \item \textbf{Seguridad en APIs:} política de desarrollo de \emph{microservicios} y
        \emph{Software Supply Chain (API) Management}. Cuantos más servicios/APIs, más
        ``cajitas'' atacables; si cada una está a cargo de un equipo distinto, puede no haber
        una política de seguridad homogénea.
  \item \textbf{RASP (Runtime Application Self Protection):} pueden estar \emph{mal
        implementados}. Son herramientas que controlan los binarios del sistema para que no
        cambien durante la ejecución; si un malware intenta modificar o reemplazar un binario,
        lo detectan y pueden alertar o tirar abajo el sistema.
  \item \textbf{Verificar Federadores:} vulnerabilidades conocidas sobre protocolos. Se
        implementan protocolos como \emph{OAuth} (p.\ ej.\ login con Google); al ser
        protocolos abiertos y antiguos, vienen siendo atacados desde su origen y una mala
        implementación los vuelve vulnerables.
  \item \textbf{Software EOL (End Of Life):} software fuera de soporte, sin actualizaciones de
        seguridad o desactualizado. Es ``muy jugoso'' a la hora de pentestear.
\end{itemize}

\subsubsection*{Examinando mecanismos y comparando con otros sistemas}
\begin{itemize}[leftmargin=1.6em,itemsep=2pt]
  \item Los \textbf{mecanismos} pueden estar mal implementados, el ambiente donde corren puede
        introducir errores, o pueden directamente no ser seguros.
  \item \textbf{Comparar con otros sistemas:} \emph{sistemas parecidos tienen problemas
        parecidos}.
\end{itemize}

\begin{defbox}{Resultado de la etapa de hipótesis}
Una \textbf{lista de posibles vulnerabilidades} que asumimos que tiene el sistema.
\end{defbox}

\begin{warnbox}{Anécdota: Supply Chain y NPM}
El \textbf{Software Supply Chain} es hoy \emph{top 3} del OWASP Top Ten para aplicaciones web
(``cuando yo empecé a trabajar ni se hablaba''). Abarca \textbf{dependencias vulnerables},
infiltraciones en flujos de \emph{CI/CD} (plugins/dependencias que ejecutan chequeos) y
\emph{dependency confusion}: una dependencia que parece genuina pero es otra, terminando en
la descarga de malware. En \textbf{NPM} el problema se agrava por las \emph{dependencias
transitivas}: NPM baja \emph{todo el árbol} de dependencias, así que si una dependencia muy
usada queda afectada, tu aplicación es vulnerable ``de por sí''. Han ido apareciendo otros
gestores que evitan bajar lo innecesario.
\end{warnbox}

\subsection{Etapa 3 — Prueba}

\subsubsection*{Priorizar}
Priorizar la lista de posibles vulnerabilidades. \emph{Depende del objetivo de la prueba};
por lo general se prioriza por \textbf{nivel de acceso requerido}.

\subsubsection*{Cómo probar la existencia de la vulnerabilidad}
\begin{itemize}[leftmargin=1.6em,itemsep=2pt]
  \item \textbf{La mejor manera:} como resultado de analizar la \emph{documentación} o el
        \emph{comportamiento}.
  \item \textbf{Otra manera:} \textbf{intentar explotarla}. Es el \emph{último recurso} y el
        \emph{menos eficiente}, porque hay que lograr un exploit como lo haría un atacante.
        \textbf{Excepción: \emph{Pivoting}} (explicado más adelante).
\end{itemize}

\begin{notebox}{Demostrar = explotar (con cuidado)}
Hoy se asume que, para \emph{demostrar} una vulnerabilidad, hay que tratar de
\textbf{explotarla} (es lo que da peso al hallazgo: ``probé la existencia explotándola''). Es
el origen de las \emph{PoC (Proof of Concept)} en pentesting. Pero hay que tener mucho
cuidado, sobre todo en sistemas productivos.
\end{notebox}

\subsubsection*{Diseñar el test — lo menos intrusivo posible}
Algunas vulnerabilidades pueden \textbf{denegar el servicio} (DoS) y atentar contra la tríada
\textbf{CIA} (Confidentiality, Integrity, Availability). \emph{Proceso normal:}
\begin{enumerate}[leftmargin=1.8em,itemsep=2pt]
  \item \textbf{Resguardar los datos} de todo el sistema (backup).
  \item \textbf{Documentar} los requerimientos para detectar la vulnerabilidad.
  \item \textbf{Intentar detectar} la vulnerabilidad.
\end{enumerate}

\begin{warnbox}{La prueba DEBE ser REPETIBLE y CONCLUSIVA}
Si funciona una sola vez y no se puede replicar (algo esporádico), \textbf{no concluye} que
haya vulnerabilidad. \emph{Repetible} significa que si el pentester puede hacerlo,
naturalmente un atacante también podrá.

\textbf{Anécdotas sobre el peligro de explotar:}
\begin{itemize}[leftmargin=1.4em,itemsep=2pt]
  \item \textbf{Chernóbil:} era una \emph{prueba controlada} (ver cómo respondía el núcleo en
        un escenario sin energía) y terminó en desastre. ``Controlado'' no garantiza seguridad.
  \item \textbf{Nmap:} programa de Linux (aparece en \emph{Matrix}: Trinity lo corre) que
        escanea interfaces y puertos y arma la \textbf{topología de la red} (un \emph{discovery}
        de la red). El profesor vio a gente \textbf{tirar abajo la red de la oficina} corriendo
        un Nmap, afectando la disponibilidad.
\end{itemize}
\end{warnbox}

\subsection{Pivoting}

\begin{defbox}{Pivoting}
La existencia de una vulnerabilidad probablemente \textbf{provea más información} (ejemplo:
acceso al Sistema Operativo). En ese caso \textbf{se debe volver a la fase de recolección de
información}. \emph{Pivoting} es el \textbf{uso de un punto de acceso comprometido para
continuar un ataque}.

\emph{Ejemplo:} se compromete el \textbf{firewall} y desde allí se continúa atacando la
\textbf{red interna}.
\end{defbox}

Esta es la \textbf{excepción} a ``explotar es el último recurso'': encadenar una
vulnerabilidad con otra hace ganar cada vez más posibilidades y ofrecer más ``daño''. El
diagrama recrea la idea del firewall comprometido como trampolín:

\begin{center}
\begin{tikzpicture}[
  >=Stealth, node distance=10mm,
  ent/.style={draw=black!60, rounded corners=2pt, minimum height=0.9cm, align=center, font=\small}]
  \node[ent, fill=attackRed!15, draw=attackRed] (att) {Atacante /\\Pentester};
  \node[ent, fill=itbaCyan!40, right=18mm of att] (fw) {Firewall\\\footnotesize(comprometido)};
  \node[ent, fill=itbaBlue!35, right=18mm of fw] (net) {Red interna\\\footnotesize(servidores)};
  \draw[->, attackRed, line width=1pt] (att) -- node[above, font=\scriptsize] {1. compromete} (fw);
  \draw[->, attackRed, line width=1pt] (fw)  -- node[above, font=\scriptsize] {2. pivotea} (net);
  % nuevo ciclo de recolección
  \draw[->, dashed, itbaCyanDark] (net.south) .. controls +(0,-1.0) and +(0,-1.0) .. (fw.south)
       node[midway, below, font=\scriptsize\itshape, text=itbaCyanDark] {vuelve a recolectar información};
\end{tikzpicture}
\end{center}

\subsection{Etapa 4 — Generalización}

A medida que las pruebas resultan exitosas, \textbf{emergen patrones}: ``esto se está dando
muy seguido, tiene que haber algo por acá''. La idea es \emph{conectar puntos}.

\begin{warnbox}{Dos vulnerabilidades combinadas = problema grave}
A veces dos vulnerabilidades combinadas constituyen un problema grave. \emph{Ejemplo de la
slide:}
\begin{itemize}[leftmargin=1.4em,itemsep=1pt]
  \item Una \textbf{cuenta de invitado} habilitada permite \textbf{conexión remota}.
  \item Un \textbf{buffer overflow local} permite \textbf{derechos de administrador}.
\end{itemize}
$\Rightarrow$ A través de una cuenta de invitado se pueden lograr \textbf{permisos de
administrador}. Una vulnerabilidad puede no ser crítica por sí sola, pero \emph{combinada} con
otra causar un daño importante.
\end{warnbox}

\begin{notebox}{La ``mano'' del pentester}
\prof{Constituye la parte más importante de la prueba.} Es donde el pentester identifica
cosas, las combina y llega a conclusiones más avanzadas. Su función es encontrar la mayor
cantidad de vulnerabilidades y ``fingir'' el mayor daño (entre comillas, no es un daño real)
para darle ese \emph{insight} a la empresa y que pueda corregir toda esa superficie de ataque.
\end{notebox}

\subsection{Etapa 5 — Eliminación (opcional)}

Por lo general \textbf{sólo se incluyen recomendaciones} (no es obligatorio, pero se suele
hacer): quien ejecuta la prueba \emph{no suele ser} el diseñador/desarrollador, así que el
dueño del sistema termina decidiendo qué hacer.

Es importante que queden claros el \textbf{contexto, los detalles y el mecanismo de
explotación}:
\begin{itemize}[leftmargin=1.6em,itemsep=2pt]
  \item Para poder \textbf{corregir} el sistema.
  \item Para poder \textbf{impedir / monitorear} mientras tanto.
  \item Para \textbf{verificar si fue explotado} y, además, poder \emph{reproducir} la prueba
        luego del fix y comprobar que ya no es vulnerable.
\end{itemize}

\begin{notebox}{El reporte de pentesting}
Todo termina en un \textbf{reporte} donde se detallan las vulnerabilidades encontradas y cómo
se propone eliminarlas (como \emph{recomendaciones}). A veces la remediación es tan simple
como actualizar la versión de una dependencia; otras veces es más abstracta (``manejar bien la
autorización''). Hay gente dedicada a documentar la remediación de los \textbf{CVE}.
\end{notebox}

%=====================================================================
\section{Ejemplo 1 — Michigan Terminal System (MTS)}
%=====================================================================

\begin{slidebox}
\centering Esta prueba de penetración contaba con la \textbf{aprobación de los
administradores}, que fueron partícipes de la prueba.
\end{slidebox}

\medskip
El \textbf{Michigan Terminal System (MTS)} es un sistema operativo que corre en
\textbf{mainframes IBM 360/370} (computadoras enormes, en general con aplicaciones bancarias).
\begin{itemize}[leftmargin=1.6em,itemsep=2pt]
  \item \textbf{Objetivo de la prueba:} obtener acceso a las \textbf{estructuras de control}
        que no debería permitirse.
  \item \textbf{Nivel inicial:} cuenta autorizada (\textbf{nivel 3}, atacante con acceso al
        sistema).
\end{itemize}

\subsection{Paso 1 — Recolección de información}

Se aprenden los detalles del \textbf{mecanismo de control}: al correr un programa, la memoria
se divide en \textbf{segmentos}.

\begin{center}
\begin{tikzpicture}[
  seg/.style={draw=black!60, minimum width=10cm, minimum height=0.95cm, align=left,
              font=\small, inner xsep=10pt}]
  \node[seg, fill=attackRed!12] (s04) {\textbf{Segmentos 0--4:} supervisor, programas de sistema y estado
        \quad\textcolor{attackRed}{\footnotesize(protegidos por hardware)}};
  \node[seg, fill=itbaCyan!35, below=1mm of s04] (s5) {\textbf{Segmento 5:} área de trabajo del sistema,
        info.\ del proceso (incl.\ nivel de privilegio)};
  \node[below=0mm of s5.south east, anchor=north east, font=\scriptsize\itshape, text=itbaCyanDark]
        {el proceso NO puede alterar esta información};
  \node[seg, fill=practGreenL, below=6mm of s5] (s6) {\textbf{Segmentos 6+:} información del proceso
        \quad\textcolor{practGreen}{\footnotesize(el proceso modifica libremente)}};
\end{tikzpicture}
\end{center}

\subsubsection*{El segmento 5 (el codiciado)}
Está protegido por el \textbf{sistema de protección de memoria virtual}:
\begin{itemize}[leftmargin=1.6em,itemsep=2pt]
  \item \textbf{Modo sistema (kernel):} el proceso puede acceder, modificar datos y ejecutar
        funciones privilegiadas.
  \item \textbf{Modo usuario:} el segmento 5 \emph{no se mapea} al espacio de direcciones del
        usuario.
\end{itemize}
Un proceso corre en \textbf{modo usuario}; desde ahí realiza \textbf{system calls}, que
corren en \textbf{modo sistema} (de forma controlada).

\subsubsection*{Ejecución de un system call}
\begin{enumerate}[leftmargin=1.8em,itemsep=2pt]
  \item El código del system call \textbf{revisa los parámetros}.
  \item Especialmente, revisa que los \textbf{punteros} pertenezcan a zonas \emph{accesibles
        por el proceso}.
  \item Los parámetros se construyen como una \textbf{lista de direcciones en el segmento de
        usuario}.
  \item La lista de parámetros se pasa como una \textbf{dirección en un registro}.
\end{enumerate}

\subsection{Paso 2 — Hipótesis}

\begin{defbox}{La hipótesis}
\textbf{¿Qué ocurriría si la dirección de un parámetro apunta a la lista misma de
parámetros?} (un puntero que apunta a sí mismo).

Sin controles extra, sería posible que algún system call \textbf{escriba la dirección del
parámetro \emph{después} de haber pasado la validación}. Si se puede controlar qué valor
escribir, técnicamente sería posible \textbf{leer o escribir cualquier zona del segmento 5}
(el segmento que no se debería poder tocar). Es un \emph{TOCTOU}: se valida y luego se altera.
\end{defbox}

\subsection{Paso 3 — Prueba}

Buscar un system call que: \textbf{(a)} use la convención de llamada estudiada, \textbf{(b)}
tenga al menos dos parámetros y altere uno, y \textbf{(c)} permita que ese parámetro se altere
con cualquier valor.

\begin{notebox}{El candidato elegido: \texttt{line input}}
\texttt{line input} lee una línea de texto y \textbf{retorna dos valores}: el \emph{número de
línea} y la \emph{longitud de línea}.

\textbf{Setup del exploit:} hacer que la \emph{dirección} donde se almacenará el
\textbf{número de línea} sea la de la \textbf{longitud de línea}.
\end{notebox}

\subsubsection*{Ejecución del exploit}
\begin{enumerate}[leftmargin=1.8em,itemsep=2pt]
  \item El system call \textbf{valida los parámetros}: todos pertenecen al segmento del
        proceso $\Rightarrow$ pasa la validación.
  \item Ejecuta el método \texttt{read input line}: la \textbf{longitud} de la línea ingresada
        es el \emph{valor a escribir}.
  \item El \textbf{número de línea} se guarda según la lista de parámetros; se logra que sea
        una \textbf{dirección del segmento 5}. Esto \textbf{reescribe la dirección del otro
        parámetro}.
  \item La \textbf{longitud} se guarda según la lista de parámetros $\Rightarrow$ esto
        \textbf{escribe el valor en el segmento 5}.
\end{enumerate}

\begin{center}
\begin{tikzpicture}[
  >=Stealth, node distance=7mm, font=\footnotesize,
  stp/.style={draw=black!60, rounded corners=2pt, fill=itbaBlue!25,
               minimum width=3.0cm, minimum height=1.0cm, align=center}]
  \node[stp] (a) {Validación OK\\\scriptsize(todo en seg.\ del proceso)};
  \node[stp, right=of a] (b) {\texttt{read input line}\\\scriptsize longitud = valor a escribir};
  \node[stp, right=of b] (c) {Nro.\ de línea\\\scriptsize $\to$ dir.\ del seg.\ 5};
  \node[stp, right=of c, fill=attackRed!15, draw=attackRed] (d) {Escritura en\\\textbf{segmento 5}};
  \draw[->] (a)--(b); \draw[->] (b)--(c); \draw[->] (c)--(d);
\end{tikzpicture}
\end{center}

\subsection{Paso 4 — Generalización}

\begin{itemize}[leftmargin=1.6em,itemsep=2pt]
  \item \textbf{No} se puede escribir en los segmentos 0--4 (protegidos por hardware).
  \item \textbf{Pero} en el segmento 5, el nivel de privilegio determina si pueden hacerse
        \textbf{llamadas de nivel supervisor} (en lugar de system calls).
  \item Y una función de nivel supervisor \textbf{apaga la protección por hardware}.
\end{itemize}

\begin{warnbox}{Conclusión del MTS}
Esta vulnerabilidad permite que un atacante \textbf{modifique cualquier dirección de cualquier
segmento}, \textbf{controlando completamente la computadora}. (El profesor aclara que era otra
época, sin los conceptos fuertes de seguridad de hoy; este ejemplo proviene del material de
Pablo.)
\end{warnbox}

%=====================================================================
\section{Ejemplo 2 — Ataque externo (ingeniería social)}
%=====================================================================

\begin{itemize}[leftmargin=1.6em,itemsep=2pt]
  \item \textbf{Objetivo:} determinar si las medidas de seguridad de una empresa eran
        efectivas para evitar que un atacante ingrese a un sistema.
  \item \textbf{Enfoque:} el test se enfoca en \textbf{políticas y procedimientos}, tanto
        técnicos como \textbf{no técnicos}.
\end{itemize}

\begin{warnbox}{Premisa: la ingeniería social}
Este pentest está fuertemente apalancado en \factor{ingeniería social}, que asume que es
\textbf{mucho más fácil atacar y vulnerar a las personas que a los sistemas}. Se aprovecha la
\emph{urgencia} y la \emph{confianza}: ``estoy apurado, lo necesito para la auditoría'' hace
que la gente entregue información ``en base a urgencia''.
\end{warnbox}

\subsection{Paso 1 — Recolección de información}

\begin{itemize}[leftmargin=1.6em,itemsep=2pt]
  \item \textbf{Búsqueda en internet:} nombres de empleados y directores; teléfono de la
        sucursal local. De ahí se consiguió el \textbf{reporte anual público} (la empresa
        cotizaba en bolsa).
  \item \textbf{Reconstrucción de estructura:} partes del \textbf{organigrama}, nombres de
        proyectos y personas que trabajan en cada uno.
\end{itemize}

\subsubsection*{Obtención del listado de empleados (encadenamiento de pretextos)}
\begin{enumerate}[leftmargin=1.8em,itemsep=2pt]
  \item El tester se hizo pasar por un \textbf{nuevo empleado} y aprendió que para enviar
        cualquier documento se necesita: \textbf{número de empleado} y \textbf{centro de
        costos}.
  \item Llamó a la \textbf{secretaria del director} (sobre quien más información había
        recabado):
  \begin{itemize}[leftmargin=1.4em,itemsep=1pt]
    \item Primero, haciéndose pasar por un empleado, consiguió el \textbf{número de empleado}
          del director.
    \item Luego, haciéndose pasar por un \textbf{auditor}, obtuvo el \textbf{centro de costos}.
  \end{itemize}
  \item Con esos datos, \textbf{solicitó enviar el listado} de empleados a una consultora
        externa.
\end{enumerate}

\subsection{Paso 2 — Hipótesis}

\begin{defbox}{Hipótesis}
Los \textbf{empleados nuevos no conocen todos los procesos y controles} de la compañía. Por lo
tanto, es posible que un nuevo empleado \textbf{brinde información sensible}.
\end{defbox}

\subsection{Paso 3 — Prueba}

\begin{enumerate}[leftmargin=1.8em,itemsep=3pt]
  \item El tester llamó a \textbf{RRHH} impersonando a la \textbf{secretaria de un director}:
        se quejó de que no le habían informado los nuevos ingresos y los pidió. Obtuvo los
        \textbf{nombres de los nuevos empleados}.
  \item El tester llamó a los \textbf{nuevos empleados} haciéndose pasar por el
        \textbf{operador del centro de cómputos} y les brindó un ``entrenamiento de seguridad
        por teléfono''. Durante el entrenamiento obtuvo: \textbf{tipos de sistemas usados,
        número de usuarios, logins y claves}.
\end{enumerate}

\begin{notebox}{El truco para sacar la clave}
El nuevo empleado confía: alguien con su número lo llamó y trabaja para la empresa. En medio
del ``entrenamiento'', el tester critica la contraseña (``es muy débil, no usaste símbolos ni
mayúsculas, no cumple la política''), se ofrece a cambiarla y pide: ``decime cuál es y yo la
cambio''. Así obtiene la clave. \emph{Resultado:} se probó que los empleados nuevos no estaban
al tanto de los procesos/controles y brindaban información sensible. La hipótesis se cumplió.
\end{notebox}

\begin{center}
\begin{tikzpicture}[
  >=Stealth, font=\footnotesize, node distance=6mm,
  role/.style={draw=black!60, rounded corners=2pt, fill=itbaBlue!25,
               minimum width=2.5cm, minimum height=0.95cm, align=center}]
  \node[role] (t1) {Falso empleado\\\scriptsize+ auditor};
  \node[role, right=10mm of t1] (sec) {Secretaria\\del director};
  \node[role, below=8mm of t1] (rrhh) {RRHH};
  \node[role, right=10mm of rrhh] (new) {Nuevos\\empleados};
  \node[role, right=10mm of new, fill=attackRed!15, draw=attackRed] (creds) {Logins\\y claves};
  \draw[->] (t1) -- node[above,font=\scriptsize,align=center]{nro.\ emp.\ +\\centro costos} (sec);
  \draw[->] (t1) -- node[left,font=\scriptsize,align=center]{impersona\\secretaria} (rrhh);
  \draw[->] (rrhh) -- node[above,font=\scriptsize]{nombres} (new);
  \draw[->] (new) -- node[above,font=\scriptsize]{``training''} (creds);
\end{tikzpicture}
\end{center}

%=====================================================================
\section{Para discutir: ¿cuán válidos y de qué calidad son los pentests?}
%=====================================================================

\subsection{¿Cuán válidos son los tests de penetración?}
\begin{itemize}[leftmargin=1.6em,itemsep=2pt]
  \item \textbf{No sustituyen} una buena especificación, diseño, implementación y pruebas
        (unitarias, de integración, cobertura, principios de diseño de seguridad).
  \item Son una técnica \textbf{importante para probar un sistema luego de ser instalado}.
        \emph{Idealmente no sería necesario; en la práctica, sí.}
  \item Encuentran \textbf{problemas introducidos por la interacción} del sistema con los
        usuarios y el ambiente, que suelen quedar \emph{fuera} del análisis y las pruebas
        normales. \emph{``Siempre algo aparece.''}
\end{itemize}

\subsection{¿Qué determina la calidad de una prueba de penetración?}
\begin{itemize}[leftmargin=1.6em,itemsep=2pt]
  \item La metodología de hipótesis de fallas \textbf{depende de la capacidad de los testers}
        para formular hipótesis (su ``visión'').
  \item \textbf{No es sistemática:} no provee una forma sistemática de revisar un sistema (a
        diferencia de la verificación formal).
  \item Los resultados de un test sirven sólo \textbf{marginalmente} para otros:
        \emph{cada test es un mundo aparte}, caracterizado por el sistema, la gente y muchos
        factores. Hay herramientas que automatizan \emph{algunos} aspectos, pero no todos.
\end{itemize}

\begin{warnbox}{Un pentest sirve para ese sistema y ese momento}
Un pentest sirve \textbf{para ese sistema en ese momento}: un mes después la API, el
comportamiento o las versiones cambiaron y la gran mayoría de las pruebas pudo quedar
\emph{obsoleta}. Por eso no se puede generalizar como sí se hace con la verificación formal.
\end{warnbox}

%=====================================================================
\section{Síntesis}
%=====================================================================

\begin{center}
\renewcommand{\arraystretch}{1.25}
\begin{tabular}{>{\raggedright\arraybackslash}p{3.4cm} >{\raggedright\arraybackslash}p{10.6cm}}
\toprule
\textbf{Concepto} & \textbf{Idea clave} \\
\midrule
Verificación formal & Prueba matemática; prueba \emph{ausencia} de vulnerabilidades (con todos los factores); costosa e impráctica. \\
Pentesting & Prueba \emph{existencia}, nunca ausencia; económico; ethical hacking; ataca el sistema como un todo. \\
Hipótesis de Falla & 5 etapas: Recolección → Hipótesis → Prueba → Generalización → Eliminación (opcional). \\
Recolección & Modelar el sistema, discrepancias, privilegios/cuentas, ambiente y red. \\
Hipótesis & Inconsistencias políticas/mecanismos; APIs/Supply Chain, RASP, federadores (OAuth), EOL. \\
Prueba & Priorizar por acceso; demostrar explotando (último recurso); test \emph{repetible y conclusivo}; backup + documentar. \\
Pivoting & Usar un acceso comprometido para seguir atacando; volver a recolectar. \\
Generalización & Conectar patrones; combinar vulnerabilidades (lo más importante). \\
Eliminación & Recomendaciones; contexto/detalle/mecanismo para corregir, monitorear y verificar. \\
\bottomrule
\end{tabular}
\end{center}

%=====================================================================
\section*{Lectura recomendada}
\addcontentsline{toc}{section}{Lectura recomendada}
%=====================================================================
\begin{slidebox}
\centering
\textbf{Capítulo 23 (secciones 1--2)} --- \emph{Computer Security: Art and Science},
Matt Bishop.\\[4pt]
\textbf{OSSTMM} --- \emph{Open Source Security Testing Methodology Manual},
\url{http://www.isecom.org/osstmm/}
\end{slidebox}

\end{document}
